
Da bisher keine praktische Implementierung der bitweisen
Encodierung existiert, besteht das erste Ziel dieser Arbeit
darin, zu zeigen, dass eine bitweise Übertragung möglich
ist und welche Grenzen sie hat. Dieser experimentelle
Abschnitt gibt erste Einblicke in die
Betrachtung der Fehlerwahrscheinlichkeiten der beiden
Encodierungen sowie in ihren Einfluss auf die Laufzeit.
Die danach ermittelte
theoretische Grenze betrachtet nicht die Korrektheit
einer Übertragung und geht davon aus, dass Übertragungen
durch den \ac{CbCC} immer korrekte Ergebnisse liefern.

\subsection{Threat-Model}
Um die Möglichkeiten der Optimierung des \ac{CbCC}
einzugrenzen, wird das fol\-gende Threat-Model für
\acp{TEA} betrachtet:

\paragraph{Der Angreifer hat}
\begin{itemize}
    \item Code-Execution auf dem
    physischen Zielsystem. (in Sandbox/VM oder als Prozess)
    \item Zugriff auf Zeitmessungen mit
    mindestens einer Millisekunde Auflösung.
    \item einen Angriff, mit dem er während eines transient
    Fensters Zugriff auf ein gesuchtes Secret hat.
    \item die Möglichkeit, während des transient Fensters eigenen
    Code auszuführen.
    \item Zugriff auf Methoden, um Cache-Zustände mindestens durch
    das Laden von Adressen zu ändern.
    \item einen geteilten Speicherbereich mit dem Zielkontext.
    (z.B. könnte der Zielkontext Zugriff auf physischen Speicher haben)
\end{itemize}

\paragraph{Verschiedene  Angriffe} aus der Vergangenheit,
die für dieses Threat Model infrage kommen sind
\cite{lipp_meltdown_2020}, \cite{maisuradze_ret2spec_2018},
\cite{schwarz_zombieload_2019} Variante 2,
\cite{kim_slap_2025}, \cite{kim_flop_2025},
\cite{kim_ileakage_2023}, \cite{hertogh_rain_nodate}.
Hierbei sei zu beachten, dass sich der Anwendungsbereich bei
fehlendem geteilten Speicherbereich leicht erweitern lässt,
da die grundlegenden Konzepte auch für \acp{CbCC} gelten,
die auf Eviction-Sets basieren, wie z.B. Prime+Probe
(\cite{liu_last-level_2015}). Primär ändert sich der Aufwand
durch das Finden von Eviction-Sets, was bereits zu einem
Aufwand von unter 100 Millisekunden optimiert wurde (\cite[3768]{kessous_pruneplumtree_2024}),
und das Füllen der Cache-Sets während der Setup-
beziehungsweise Prime-Phase.
